\section{Annotation Review and Failure Taxonomy}
\label{sec:audit}

The automatic label audit checks pairing, normalized coordinates, zero-area
boxes, frame excursions, empty labels, annotation density, normalized box
area, and possible image duplication.  It does not prove semantic correctness.
A separate review sheet therefore samples boxes across crop, source, box-size
stratum, image boundary, and annotation-density stratum.  Reviewers mark
correct, loose, tight, duplicate, wrong-object, or ambiguous.  Because the
frozen artifacts do not contain adjudicated review sheets, this supplement
makes no semantic annotation-acceptance claim.  This distinction prevents an
automated format audit from being misreported as human semantic validation.

\paragraph{Failure taxonomy.}
Qualitative analysis uses balanced denominators.  Detector failures comprise
misses, background false positives, duplicates, and boundary truncation; VLM
failures comprise task errors, malformed schema, unsupported claims, and repeat
inconsistency.  Gate outcomes distinguish correct/false acceptance from
correct/unnecessary abstention.  Displayed cases identify crop, protocol,
ground truth, prediction, selector, budget, and gate decision.

Hard-negative mining may be evaluated as a detector-training ablation only
when the mined negatives, held-out split, and before/after checkpoints are
available under the same source-safe protocol.  Results from a separate
pipeline or a screenshot are not merged into the reported experiment.
